\section{\MethodName: Chunk-based Outline-guided Patent Generation}

Since LLMs cannot yet generate sufficiently long documents in a single call, we instead choose to generate patents in chunks. 
\Cref{fig:generation_approach} summarizes the approach.
We chunk the outline, select the most relevant parts of the paper (\textit{paper context})
depending on the chunks' outline bullet points, and prompt an LLM to generate the patent text for that chunk.
The outlines of previous chunks are included for global document context.
Finally, we concatenate the generated outputs and apply only a lightweight post-processing to remove duplicate headings at chunk boundaries.
This framework enables generating a long document in parallel, with customized context per chunk and controllable length.


\begin{table*}[t]
    \centering
    \small
    \resizebox{\textwidth}{!}{
      \setlength{\tabcolsep}{4pt}
      \begin{tabular}{lccccccccccc}
        \toprule
                              & \multirow{4}{*}{Tokens}
                              & \multicolumn{2}{c}{\multirow{3}{*}{Text Sim $\uparrow$}}
                              & \multicolumn{3}{c}{Content-Level Metrics (SCALE) $\uparrow$}
                              & \multicolumn{2}{c}{\multirow{3}{*}{Language $\uparrow$}}
                              & \multicolumn{2}{c}{\multirow{3}{*}{Repetitions}}
        \\
        \cmidrule(lr){5-7}
                              &
                              &
                              &
                              & Coverage
                              & \multicolumn{2}{c}{Factuality}
                              &
                              &
                              &
                              &
                              &
        \\
        \cmidrule(lr){3-4}\cmidrule(lr){6-7}\cmidrule(lr){8-9}\cmidrule(lr){10-11}
                              &
                              & BS
                              & R-L
                              & \textit{Gen}\textrightarrow{}\textit{Ref}
                              & \textit{Ref}\textrightarrow{}\textit{Gen}
                              & \textit{Ref}+\textit{Pap}\textrightarrow{}\textit{Gen}
                              & Style
                              & DS
                              & RR
                              & RR>80
        \\
        \midrule
    
        \multicolumn{12}{c}{\textbf{Heuristic Baselines / Skylines}}\vspace{.5em}                                                                                                      \\
        Reference Patent
                              & \textit{18.1k (100\%)}
                              & \textit{100}
                              & \textit{100}
                              & \textit{88.6 $\pm$ .15}
                              & \textit{88.5 $\pm$ .18}
                              & \textit{88.7 $\pm$ .18}
                              & \textit{100}
                              & \textit{100}
                              & 14.4
                              & 0.2
        \\
        Similar Patent
                              & 30.1k (166.0\%)
                              & 65.8
                              & 33.7
                              & 29.4 $\pm$ .37
                              & 27.6 $\pm$ .32
                              & 27.6 $\pm$ .36
                              & 75.3
                              & 98.3
                              & 12.3
                              & 0.1
        \\
        Outline
                              & 1.4k (7.9\%)
                              & 56.5
                              & 10.1
                              & 39.7 $\pm$ .72
                              & 61.6 $\pm$ .14
                              & 61.9 $\pm$ .15
                              & 24.4
                              & 85.6
                              & 20.1
                              & 0.2
        \\
        Paper
                              & 8.1k (44.5\%)
                              & 69.7
                              & 39.4
                              & 44.8 $\pm$ .61
                              & 46.5 $\pm$ .40
                              & \textit{89.0 $\pm$ .13}
                              & 47.2
                              & 98.4
                              & 8.4
                              & 0.0
        \\
        \cmidrule{1-11}
        \multicolumn{12}{c}{\textbf{Single LLM-call} ($\mathcal{T}$\textit{$\{$inst=2k; pap=$\infty$; pat=$\infty\}$})}\vspace{.5em}                                                                                                          \\
        Mixtral-8x7B
                              & 3.2k (17.7\%)
                              & 66.1
                              & 23.1
                              & 38.1 $\pm$ .70
                              & \textbf{66.9 }$\pm$ .87
                              & 72.0 $\pm$ .71
                              & \textbf{\itshape 42.6}
                              & 96.9
                              & 17.9
                              & 0.6
        \\
        Qwen2-72B
                              & 2.8k (15.6\%)
                              & \textbf{\itshape 66.6}
                              & 21.3
                              & \textbf{\itshape 40.3 }$\pm$ .45
                              & 65.8 $\pm$ .47
                              & 72.4 $\pm$ .43
                              & 39.6
                              & 97.3
                              & 9.5
                              & 0.5
        \\
        $\quad$w/o Paper
                              & 3.5k (19.3\%)
                              & 66.1
                              & \textbf{\itshape 23.2}
                              & 38.9 $\pm$ .29
                              & 65.9 $\pm$ .54
                              & 66.6 $\pm$ .38
                              & 37.1
                              & 96.6
                              & 9.8
                              & 0.6
        \\
        $\quad$w/o Outline
                              & 2.0k (11.1\%)
                              & 64.2
                              & 16.9
                              & 34.9 $\pm$ .40
                              & 56.7 $\pm$ .31
                              & \textbf{75.3 }$\pm$ .23
                              & 39.8
                              & \textbf{\textit{97.4}}
                              & 9.3
                              & 0.1
        \\
    
        \cmidrule{1-11}
        \multicolumn{12}{c}{\textbf{\MethodName} ($\mathcal{T}$\textit{$\{$inst=2k; pap=3k; pat=2k$\}$})}\vspace{.5em}                                                                                \\
        Llama-3 8B
                              & 9.6k (53.1\%)
                              & 68.7
                              & 41.4
                              & 40.3 $\pm$ .21
                              & 60.8 $\pm$ .49
                              & 65.7 $\pm$ .55
                              & 43.2
                              & 97.0
                              & 27.0
                              & 4.4
        \\
        Llama-3 8B SFT
                              & 27.5k (151.5\%)
                              & 70.4
                              & 42.8
                              & 42.0 $\pm$ .25
                              & 49.3 $\pm$ .39
                              & 52.1 $\pm$ .39
                              & 59.4
                              & 98.0
                              & 53.7
                              & 29.3
        \\
        $\quad$w/ rep. removal
                              & 17.3k (95.4\%)
                              & \textbf{\itshape 71.2}
                              & \textbf{\itshape 49.6}
                              & 42.1 $\pm$ .48
                              & 49.6 $\pm$ .38
                              & 53.4 $\pm$ .59
                              & \textbf{64.7}
                              & \textbf{98.5}
                              & 38.4
                              & 8.5
        \\
        Mixtral-8x7B
                              & 5.6k (31.0\%)
                              & 69.1
                              & 35.5
                              & 41.8 $\pm$ .61
                              & 62.3 $\pm$ .31
                              & 68.5 $\pm$ .30
                              & 49.3
                              & 97.5
                              & 14.5
                              & 0.4
        \\
        Llama-3 70B
                              & 6.1k (33.9\%)
                              & 70.2
                              & 39.0
                              & 42.7 $\pm$ .60
                              & \textbf{\itshape 64.5 }$\pm$ .47
                              & \textbf{\itshape 68.6 }$\pm$ .58
                              & 49.5
                              & 97.4
                              & 17.5
                              & 0.2
        \\
        Qwen2-72B
                              & 8.1k (44.8\%)
                              & 70.2
                              & 41.3
                              & \textbf{\itshape 44.1 }$\pm$ .32
                              & 62.5 $\pm$ .44
                              & 67.9 $\pm$ .35
                              & 47.5
                              & 97.3
                              & 12.8
                              & 1.2
        \\
        \cmidrule{1-11}
        \multicolumn{12}{c}{\textbf{\MethodName} ($\mathcal{T}$\textit{$\{$inst=2k; pap=3k; pat=400$\}$})}\vspace{.5em}                                                                                    \\
        Qwen2-72B
                              & 18.1k (100\%)
                              & \textbf{71.7}
                              & \textbf{50.8}
                              & \textbf{46.8 }$\pm$ .31
                              & 59.7 $\pm$ .63
                              & 65.3 $\pm$ .44
                              & 47.8
                              & 97.1
                              & 10.0
                              & 0.5
        \\
        \bottomrule
    \end{tabular}
    
    }
    \caption{
      Experimental Results. 
      \textit{Italic values} represent upper bounds that used test data in the prediction. 
      The best value per column is \textbf{bold}, the best per section \textbf{\itshape bold italics}.
      Tokens are reported as absolute and relative to the reference.
      BS = BERTScore. R-L = ROUGE-L. DS = DiscoScore. Text Sim = Standard Text Similarity Metrics.
    }
    \label{tab:results}
\end{table*}
  


\noindent\textbf{Token Allocation and Chunking. } 
In the default setting, we allocate per chunk 2k tokens for the instruction, 3k tokens for the paper context, and 2k tokens for the output patent ($\mathcal{T}$\textit{$\{$inst=2k; pap=3k; pat=2k$\}$}).
We choose this setting because preliminary experiments show that on our task, LLMs only generate up to roughly 2-3k tokens regardless of the requested amount of content,
and because this allows comparing with models that support up to 8k tokens.
The chunking procedure segments the outline into chunks depending on the token allocation.
In particular, it determines the number of outline bullet points per chunk using the reserved number of patent tokens and the average number of characters per outline bullet point\footnote{We translate between number of tokens and number of characters using average corpus statistics of patents in \DatasetName.}.
It keeps sections intact whenever possible.

\noindent\textbf{Length Control Mechanism. }
\citet{baiLongWriterUnleashing102024} find that LLMs' response lengths are constrained within a certain range and only moderately adjust to length requests in the prompt, making instructions unreliable for controlling length. 
However, as the requested length decreases, the LLM is more likely to meet it.
Hence, in \MethodName, we decrease the number of allocated patent tokens per chunk, thereby increasing the number of chunks, until the desired total length is reached.
In our experimental setting, we do this on a corpus-level, i.e., we search for a setting where the average length matches that of the references, leading to the calibrated token allocation $\mathcal{T}$\textit{$\{$inst=2k; pap=3k; pat=400$\}$}.


\noindent\textbf{Paper Context Selection. }
For each chunk $i$, the retriever selects relevant paragraphs from the paper to create the paper context $\text{c}_i$.
We use BM25 \cite{robertsonProbabilisticRelevanceFramework2009} with the chunk's outline $\text{o}_i$ as query.
We always include the abstract of the paper and all headings, adding paper paragraphs successively %
in the order of their relevance ranking until the token limit is reached.


\noindent\textbf{Patent Generation. } 
The paper context $\text{c}_i$, the current outline $\text{o}_i$ and prior outlines $\text{o}_{j<i}$ are combined to a prompt (see Appendix \ref{sec:generation_prompt_appendix}). 
The LLM generates patent chunk $\text{p}_i$, using constrained decoding to adhere to the outline's headings.
